% ============================================
% Response to Reviewer 1
% ============================================
\section*{Response to Reviewer 1}

We thank the reviewer for their constructive feedback. Below we address each point in detail.

% ============================================
\subsection*{R1-1. Absolute biomass as an alternative explanation for Dominance}

\reviewercomment{I found the Authors' explanation of coalescence generally convincing. However, I wonder if a simple, alternative explanation for many of the results could be that parent communities attain very different absolute densities. For example, if community A reaches 10x the overall density of community B, then even relatively rare species in community A might be initially more abundant than the dominant species in community B (because communities are combined in equal volumes, line 119). This might lead to the final community being dominated by species from A simply because species survive proportionally to their initial \emph{absolute} abundance. If I understand correctly, this scenario is different from the one tested in ED Fig.~5. High heterogeneity in absolute densities could potentially explain both the prevalence of Dominance outcomes and the difference in pairwise selection correlation between same- and cross-parent pairs. In principle, heterogeneity in absolute densities could also increase with nutrient concentration. To be clear, I have no reason to expect this pattern, but I wonder if the Authors could do more to rule it out. Since ODs were measured throughout (line 379), could those data be used to check this possibility?}

\responselabel
We thank the reviewer for raising this important alternative explanation. We tested the suggested hypothesis by asking whether parental-community OD$_{600}$ correlates with the parental dominance direction or with pairwise selection correlation.
If Dominance were primarily driven by unequal absolute biomass at mixing, the higher-OD parental community should win more frequently, and signed $\Delta\text{OD}_{A-B}$, the difference between the OD of parental community A and the OD of parental community B, should be positively associated with PDI\@. We observed the opposite pattern. In Base, the higher-OD parental community wins only $37\%$ of Dominance events, and the continuous relationship between signed $\Delta\text{OD}_{A-B}$ and PDI is weak. This opposite-to-prediction pattern becomes strongest in Nutr$+$, where Dominance is most frequent: the higher-OD parental community wins only $9/68 = 13.2\%$ of Dominance events ($p = 3.9 \times 10^{-10}$), and signed $\Delta\text{OD}_{A-B}$ is strongly negatively correlated with PDI ($\rho = -0.60$, $p = 4.0 \times 10^{-10}$; Response Fig.~R1-1A,B). Nutr$-$ is less informative for this directional test because both OD differences and Dominance frequency are smaller. Thus, this directional analysis argues against parental-community biomass differences as an explanation for Dominance direction.

This density-based explanation also predicts that higher-OD parental communities should show stronger within-community pairwise selection correlations. To test this, we binned parental communities into low, mid, and high OD tertiles within each medium and computed the pairwise selection correlation among species from the same parental community in each OD bin (Response Fig.~R1-1C). The cross-community correlation is not re-binned by OD and is shown as a medium-level reference. Within-community pairwise selection correlation does not show a consistent increase with parental-community OD; instead, the pattern is not consistent across media. In Nutr$+$, the slight negative OD trend mirrors the negative signed $\Delta\text{OD}_{A-B}$--PDI association in the directional analysis above, rather than the positive trend expected if higher biomass drove winning. Overall, this analysis implies that differences in parental-community biomass do not explain Dominance direction or the nutrient-dependent pattern of correlated species fates.

We added a concise note to the Results and present the corresponding biomass and richness controls in Supplementary Note~5 and Supplementary Figs.~13--15, 45. The Results now state: \mschange{``We next tested whether nutrient-dependent changes in parental-community properties, including parental-community biomass heterogeneity and initial richness effects, could explain this trend; however, Dominance outcomes were not biased toward the higher-biomass parental community, and these factors did not account for the nutrient-dependent increase in Dominance or correlated species fates (Supplementary Note~5 and Supplementary Figs.~13--15, 45).''}

\begin{figure}[H]
\centering
\includegraphics[width=\textwidth]{Fig_R1_1A_winner_loser_OD.pdf}
\caption*{\textbf{Response Fig.~R1-1A.} \textbf{Higher parental-community OD does not predict the winner in Dominance events.} Winner OD vs loser OD scatter is shown per medium; dashed line is $y=x$. Points above the line are cases where the winning parental community had higher OD than the losing parental community. The higher-OD parental community wins only $37\%$, $37\%$, and $13\%$ of Dominance events in Nutr$-$, Base, and Nutr$+$ respectively; the strongest opposite-to-prediction pattern occurs in Nutr$+$. The displayed $p$ value is from a two-sided binomial test against the null expectation that the higher-OD parental community wins $50\%$ of Dominance events.}
\end{figure}

\begin{figure}[H]
\centering
\includegraphics[width=\textwidth]{Fig_R1_1B_OD_vs_PDI.pdf}
\caption*{\textbf{Response Fig.~R1-1B.} \textbf{Continuous OD differences do not support the prediction that the higher-OD parental community wins.} Signed $\Delta\text{OD}_{A-B}$ ($\text{OD}_{A} - \text{OD}_{B}$) vs PDI is shown per medium. PDI is the Parental Dominance Index defined in the manuscript, $\text{PDI} = (2/\pi)\operatorname{atan2}(u,v)$, where $\operatorname{atan2}$ is the two-argument arctangent. Spearman $\rho$ is the rank correlation between signed $\Delta\text{OD}_{A-B}$ and PDI; the displayed $p$ value is from a two-sided Spearman rank-correlation test of the null hypothesis of no monotonic association. The relationship is weak in Base and strongly negative in Nutr$+$ ($\rho = -0.60$, $p = 4 \times 10^{-10}$), opposite the sign predicted if higher OD caused a parental community to win. Gray points are reflections included for symmetry.}
\end{figure}

\begin{figure}[H]
\centering
\includegraphics[width=\textwidth]{Fig_R1_1C_pairwise_corr_vs_OD.pdf}
\caption*{\textbf{Response Fig.~R1-1C.} \textbf{Higher parental-community OD does not produce a consistent increase in within-community pairwise selection correlation.} Columns show Nutr$-$, Base, and Nutr$+$, respectively. \textbf{Top row:} parental communities are grouped into low, mid, and high within-medium OD tertiles; smaller text below the group labels shows the OD$_{600}$ range of each tertile. Bars show within-community pairwise selection correlation with bootstrap 95\% CIs; the dashed gray line and shaded band show the medium-level mean cross-community correlation and its bootstrap 95\% CI. \textbf{Bottom row:} per-event parental-community observations plotted by OD and within-community pairwise selection correlation; vertical dotted lines mark the low/mid/high OD boundaries used in the top row. Spearman $\rho$ and $p$ values are shown in each panel; p-values are from two-sided Spearman rank-correlation tests of the null hypothesis of no monotonic association between OD and within-community pairwise selection correlation.}
\end{figure}

% ============================================
\subsection*{R1-2. Does pH mismatch predict Dominance?}

\reviewercomment{To support the hypothesis that the top-down regime emerges because dominant species strongly modify the pH of the ecosystem, the Authors show that the predictive power of the parent community pH increases in the Nutr+ condition when looking at coalescence between communities where one parent is acidic and one is alkaline (lines 261, ED Fig.~8). It seems like an even more straightforward prediction would be that Dominance should become more likely when combining parent communities with qualitatively different pH (e.g.\ alkaline and acidic) compared to cases where both parents are either acidic or alkaline. Is this the case?}

\responselabel
We thank the reviewer for suggesting this direct test. We separated two related questions: whether pH contrast increases the probability of a Dominance-class outcome, and whether pH contrast predicts which parent wins within acidic--alkaline pairs.

Acidic--alkaline pairings were more likely to produce a Dominance-class outcome than same-pH pairings in Nutr$+$: 29/32 acidic--alkaline events (90.6\%) were classified as Dominance, compared with 22/34 same-pH events (64.7\%, OR $= 5.27$, Fisher $p = 0.018$). In Base medium, acidic--alkaline pairings were not enriched for Dominance (24/41, 58.5\%) relative to same-pH pairings (23/32, 71.9\%, OR $= 0.55$, Fisher $p = 0.33$). Thus, pH contrast is associated with higher Dominance frequency in Nutr$+$ under this comparison, but not in Base, and same-pH Nutr$+$ pairings still frequently produce Dominance. This supports pH contrast as a contributor to the nutrient-rich top-down regime rather than as a complete explanation for Dominance across all coalescence events (Response Fig.~R1-2 and Supplementary Fig.~46).

This complements the winner-direction analysis already reported in the manuscript: acidic parental communities dominated 29/32 Nutr$+$ acidic--alkaline events (90.6\%), compared with 23/41 Base events (56.1\%, Extended Data Fig.~7).

We revised the overall claim in Results \S2.5 to make clear how the pH-contrast analysis bears on this question. The revised Results paragraph states, in part: \mschange{``In coalescence events between acidic (pH $<$ 6.5) and alkaline (pH $>$ 7.5) parental communities, the acidic community dominated in only 56\% of cases in Base medium ($n = 41$) and in 91\% of cases in Nutr$+$ medium ($n = 32$) (Extended Data Fig.~7). The difference between Base and Nutr$+$ proportions was significant (two-sided Fisher exact test, $p = 0.0015$), consistent with the hypothesis that high nutrients amplify metabolic activity, thereby intensifying pH modification and excluding pH-sensitive species \citep{Ratzke2018, Ratzke2020}. Together, the dominant-species and strict pH-contrast analyses suggest that Nutr$+$ Dominance may reflect a top-down regime, in which dominant taxa and associated environmental feedbacks help determine winner identity. However, this pH-based analysis addresses winner direction in high-contrast acidic--alkaline pairings and does not by itself account for the full Dominance pattern across all coalescence events. Dominance in Base is less directly explained by the measured dominant-species and pH effects tested here.''}

\begin{figure}[H]
\centering
\includegraphics[width=\textwidth]{Fig_R1_2_acidalk_per_medium.pdf}
\caption*{\textbf{Response Fig.~R1-2.} \textbf{Strict parental-community pH contrast is most informative in Nutr$+$.} Base and Nutr$+$ panels show stacked Dominance, Mixture, and Restructuring fractions for strict acidic--alkaline and strict same-pH pairings. Parental communities were classified as acidic when pH $<6.5$ and alkaline when pH $>7.5$. Intermediate-pH pairings were excluded. Statistical tests compare Dominance with pooled Mixture and Restructuring outcomes. Strict acidic--alkaline pairings were enriched for Dominance in Nutr$+$ (29/32 versus 22/34 events, OR $=5.27$, Fisher $p=0.018$), but not in Base medium (24/41 versus 23/32 events, OR $=0.55$, Fisher $p=0.33$). The Acid-Alk pairs-only panel shows winner direction within strict acidic--alkaline pairings: the acidic parental community dominated 29/32 Nutr$+$ events (90.6\%), compared with 23/41 Base events (56.1\%, Extended Data Fig.~7). Dots show individual binary winner-direction outcomes, and squares show mean $\pm$ s.e.m.}
\end{figure}

% ============================================
\subsection*{R1-3. Circularity in Fig.~5C, PDI excluding dominant species}

\reviewercomment{In Fig.~5C, it seems like there is some risk of circularity if the dominant species are included in the calculation of the PDI. Does the positive correlation under Nutr+ conditions hold up if the dominant species are excluded from this calculation?}

\responselabel
We agree that this is an important circularity check. We used two removal strategies to avoid biasing the result toward the dominant species in the coalesced community. In the first, we removed the most abundant species in the coalesced community from all three community composition vectors; in the second, stricter test, we identified the most abundant species in each parental community and removed those taxa from all three composition vectors before recalculating PDI\@. We followed the same PDI calculation and filtering criteria as in the original Fig.~5C analysis: events are included if they are classified as Mixture or Dominance before dominant species removal.

This analysis reveals that our core conclusion remains supported: dominant-species pairwise interactions remain associated with the recalculated community-level PDI in Nutr$+$, but not in Base. In Base, $R^2$ drops from $0.11$ to $0.00 / 0.07$ and the slope from $0.49$ to $-0.02 / 0.33$ (Original / coalesced-community dominant removed / parental-community dominants removed). In Nutr$+$, $R^2$ drops from $0.49$ to $0.24 / 0.18$ and the slope from $1.03$ to $0.64 / 0.56$ (Response Fig.~R1-3). Using independent, non-reflected events as the primary inference, the Spearman correlation between dominant-species pairwise score and community-level score remains significant after both removal strategies: $\rho = 0.41$ ($p = 8.2 \times 10^{-4}$) after removing the coalesced-community dominant species and $\rho = 0.42$ ($p = 1.3 \times 10^{-3}$) after removing the dominant species from both parental communities. Thus, dominant taxa contribute substantially to the Nutr$+$ signal, but the Nutr$+$ association does not disappear when those taxa are excluded. In contrast, Base does not retain a significant Spearman correlation after removal ($p = 1.00$ and $0.21$).

We added a concise note to Results \S2.5 and present the primary analysis in Supplementary Fig.~16. The Results now state: \mschange{``The Nutr$+$ predictability trend remained significant after recalculating PDI with dominant species excluded (Spearman rank-correlation tests: $p = 8.2 \times 10^{-4}$ after removing the coalesced-community dominant species and $p = 1.3 \times 10^{-3}$ after removing the two parental-community dominant species; Supplementary Fig.~16), indicating that it is not solely a dominant-species artifact in the PDI calculation.''}

\begin{figure}[H]
\centering
\includegraphics[width=\textwidth]{Fig_R1_3_per_medium_scatter.pdf}
\caption*{\textbf{Response Fig.~R1-3.} \textbf{Nutr$+$ predictability persists after dominant-species removal, arguing against a dominant-species artifact in the PDI calculation.} Per-medium scatter plots are shown with no pooling between Base and Nutr$+$. Rows: Base (MN, brown) and Nutr$+$ (HN, orange). Columns: original Fig.~5C calculation; dominant species from the coalesced community removed from all three composition vectors; and dominant species from each parental community removed from all three composition vectors before recalculating PDI\@. Colored points are independent coalescence events; gray points are reflected counterparts plotted for symmetry. The third column is the stricter test. $R^2$ and slope annotations are computed on the reflected plotting data, as in Fig.~5C, and are used as descriptive summaries. Spearman $\rho$ and $p$ are computed on the independent, non-reflected events; Spearman p-values are from two-sided rank-correlation tests of the null hypothesis of no monotonic association.}
\end{figure}

% ============================================
\subsection*{R1-4. Pool-size effects}

\reviewercomment{I wondered if the Authors could provide slightly more detail on the effect of the initial pool size. Is it similarly true that there is no significant effect of initial richness in the experimental communities? By eye, it looks like Dominance might be more likely for the 24 species pool, but it is hard to tell by squinting at Fig.~1E. It might also be interesting to see the survival ratio as a function of initial richness for the experimental communities.}

\responselabel
We did not detect a significant effect of pool size in the experiment, and we thank the reviewer for prompting this more direct check. Response Fig.~R1-4C shows Dominance frequency by initial pool size within each nutrient condition. Across nutrient conditions, the full three-category outcome distribution did not vary significantly with pool size ($\chi^2 = 2.24$, $p = 0.69$), and binary Dominance frequency was also not significantly associated with pool size within individual nutrient conditions (Dominance vs non-Dominance tests: Nutr$-$ $p = 0.39$, Base $p = 0.82$, Nutr$+$ $p = 0.27$). Thus, although initial pool size affects realized richness and ASV richness per inoculated isolate (panels A,B), these assembly-richness effects do not translate into a detectable pool-size dependence of Dominance frequency with the available sampling.

\reviewercomment{I also wondered if the Authors could look more closely at why there is no effect in the model. Perhaps by plotting the realized richness and interaction strength in parent communities as a function of initial richness?}

As a complementary model check, we performed a gLV pool-size ablation study at low, intermediate, and high values of the interaction-strength parameter $\mu$, varying the initial pool size from 4 to 48 species per parental community (100 replicate assemblies per $\mu$ and pool size). Consistent with the reviewer's intuition, reconstructing the simulated interaction matrices showed that assembly keeps within-community interaction coefficients lower than between-community coefficients across pool sizes. The outcome-level model check gives the same interpretation: at fixed $\mu$, Dominance frequency changes only modestly with pool size compared with the strong effect of $\mu$ itself (panel E), whereas the same-parental-community versus cross-parental-community pairwise selection correlation gap increases with pool size (panel F). Thus, pool size can modulate the strength of correlated species fates, but it does not fully explain the primary interaction-strength-dependent transition in Dominance frequency.

We added this analysis to Supplementary Note 5, ``Alternative models and controls for community-level selection,'' and to Supplementary Fig.~13. The new supplementary text states that pool-size variation affects assembly richness but does not explain the Dominance-frequency pattern, and that the model pool-size ablation supports the same conclusion.

\begin{figure}[H]
\centering
\includegraphics[width=\textwidth]{pool_size_analysis.pdf}
\caption*{\textbf{Response Fig.~R1-4.} \textbf{Dominance frequency is primarily associated with nutrient condition or model interaction strength, rather than initial pool size.} \textbf{A:} Realized parental richness. \textbf{B:} ASV richness per inoculated isolate, used as an assembly-retention proxy and computed as realized parental ASV richness divided by inoculated species-pool size; this ratio can exceed one because the numerator is thresholded ASV richness rather than a count restricted to the inoculated isolates. \textbf{C:} Experimental Dominance fraction by medium and initial pool size; the displayed $\chi^2$ test refers to the full three-category outcome distribution across pool sizes. \textbf{D:} Experimental pairwise selection correlation metric grouped by initial pool size. \textbf{E:} Model Dominance fraction grouped by $\mu$. \textbf{F:} Model pairwise selection correlation metric across pool sizes $4$--$48$.}
\end{figure}

% ============================================
\subsection*{R1-5. Similarity-metric robustness claim}

\reviewercomment{The Authors write (line 134) that ``This pattern of Dominance as the most frequent outcome was robust across variants of similarity metrics''. I liked the Author's approach to measuring outcomes, and I don't advocate for changing it, but this claim seems a little misleading. Two of four alternative similarity metrics gave different results, with Dominance being the \emph{least} likely outcome under Jensen--Shannon and Jaccard.}

\responselabel
We agree that our original robustness claim was too broad and could be misleading: the Dominance pattern is recovered by Euclidean distance and Bray--Curtis dissimilarity, but not under Jaccard index or Jensen--Shannon divergence (Extended Data Fig.~2). This pattern is consistent with the fact that the metrics emphasize different aspects of community composition: abundance-weighted parental-similarity metrics such as vector decomposition, Euclidean distance, and Bray--Curtis dissimilarity are more sensitive to quantitative dominance in relative abundance, whereas Jaccard reflects presence/absence retention and Jensen--Shannon evaluates divergence across the full abundance distribution. Thus, a coalesced community can be skewed toward one parental community in abundance space while still retaining taxa from both parental communities, or while redistributing subdominant taxa in ways that make it less skewed toward one parental community under Jaccard or Jensen--Shannon.

We revised Results \S2.1 to state the mixed metric result directly: \mschange{``Alternative metric analysis revealed that Dominance remained the most frequent outcome under abundance-weighted compositional metrics, including Euclidean distance and Bray--Curtis dissimilarity, but not under Jaccard index or Jensen--Shannon divergence, which emphasize species-identity retention and full-distribution divergence, respectively (Extended Data Fig.~2).''} We also revised the Extended Data Fig.~2 caption to explicitly discuss the divergent metrics: \mschange{``Dominance is recovered by abundance-weighted parental-similarity metrics, including the primary similarity-map classification, Euclidean distance, and Bray--Curtis dissimilarity, but not by Jaccard index or Jensen--Shannon divergence. This divergence reflects the different biological questions asked by each metric: abundance-weighted metrics quantify whether the coalesced community quantitatively resembles one parent more than the other, whereas Jaccard index measures presence/absence retention and Jensen--Shannon divergence measures differences across the full abundance distribution. Thus, Jaccard and Jensen--Shannon are best interpreted as complementary metric choices that emphasize different features of community composition.''}


% ============================================
\subsection*{R1-6. Gray reflected points in figures}

\reviewercomment{It took me quite a while to realize that the gray points in Figs.~1E, 4C, 5C, and 6B are simply a reflection of the colored points. Please indicate this clearly, or perhaps just remove these.}

\responselabel
We retained the reflected points and updated the captions of Figs.~1E, 4C, 5C, and 6B to state explicitly that the gray points are reflections: \mschange{``Gray points represent the symmetric counterpart of each coalescence event (community B versus A).''}

% ============================================
\subsection*{R1-7. Interaction matrix after assembly}

\reviewercomment{In Fig.~2A, it might be helpful to show the matrix of interaction coefficients after assembly of the parent communities (with visible block structure).}

\responselabel
We agree that this visualization is helpful. We added the requested post-assembly interaction-matrix analysis as Supplementary Fig.~3, where post-assembly survivor matrices are ordered by parental-community origin, and kept Fig.~2A focused on the simulation workflow. We also revised Fig.~2C and the Results text to emphasize the assembly-driven decrease in within-community interaction strength.

% ============================================
\subsection*{R1-8. Fig.~2D visualization}

\reviewercomment{Fig.~2D, I was somewhat confused by the relationship between points and squares (means?). Visually, the squares do not look like they are near the means of the points. Please clarify this for readers. Additionally, the gray horizontal bars are not explained in the main text unless I missed it.}

\responselabel
In the simulation panel, the dots show a stratified subset of per-event pairwise selection correlations for visual legibility, whereas the squares and error bars show the mean $\pm$ s.e.m.\ across all simulated coalescence events. We revised the Fig.~2D caption to specify (1) what the dots represent, (2) what the squares and error bars represent, and (3) that the gray horizontal lines are the random-selection baseline from a parental-affiliation-shuffle null model: \mschange{``Individual dots show per-event pairwise selection correlations; in the simulation panel, 100 stratified events per group are displayed for legibility. Squares with error bars show mean $\pm$ s.e.m.\ across all coalescence events ($n = 1{,}200$ per group in simulations and all valid experimental events in the experimental panel). Gray horizontal lines indicate the random selection baseline from a null model in which parental-affiliation labels are shuffled (see Supplementary Information).''}


% ============================================
\subsection*{R1-9. Emphasize ``cohesion without cooperation''}

\reviewercomment{Discussion, lines 342--345: I think it is worth emphasizing how interesting and perhaps surprising it is to find ``community-level cohesion without cooperation'' as the paper by M.\ Tikhonov has it. This finding will be counterintuitive to many readers, even though the mechanism is made clear.}

\responselabel
We agree with the reviewer that this is a central and counterintuitive implication of the study. We added a dedicated short paragraph to the Discussion that frames the result in these terms and situates this point in relation to Tikhonov's prediction and random Lotka--Volterra theory.

\mschange{``Notably, and perhaps counterintuitively, our model shows that community-level selection-like patterns can emerge without explicitly cooperative interactions between species, consistent with predictions from Tikhonov's resource-competition models \citep{Tikhonov2016} and random Lotka--Volterra theory \citep{May1972, Grilli2017}. Community-level selection is often attributed to cooperative mechanisms such as cross-feeding or division of labor. By contrast, the model suggests that assembly history and competitive interactions can be sufficient to generate correlated species fates during coalescence. Thus, cooperative and competitive mechanisms may represent distinct routes to community-level selection, and understanding how these routes interact remains an important direction for future work.''}


% ============================================
\subsection*{R1-10. Means missing in Extended Data Fig.~5C}

\reviewercomment{In ED Fig.~5C, are means missing?}

\responselabel
We thank the reviewer for pointing this out. The means are shown in Extended Data Fig.~5C as square markers with error bars, but they were not sufficiently clear in the previous version. We revised Extended Data Fig.~5 so that the mean markers and error bars are more visible, and the caption now states that individual dots show per-event correlations whereas squares with error bars show mean $\pm$ s.e.m.


% ============================================
\subsection*{R1-11. Incorrect Extended Data reference in the SI}

\reviewercomment{On p.~9 in the SI, I believe the sentence ``Simulations were performed with parental community sizes ranging from 4 to 48 species per community (Extended Data Fig.~5)'' should refer to ED Fig.~4.}

\responselabel
We thank the reviewer for catching this error and corrected the SI reference.
